\subsection{Reducing an instrument preserves the purchase and its earlier hedge history}
\label{instrument-only-reduction}

An instrument addition starts a new contribution to hedge gains. An
instrument reduction stops part of the future contribution. It does not erase
the gains already accumulated or necessarily reduce the purchase being hedged.
IFRS~9.B6.5.18 leaves measurement of the hedged item unchanged when the
instrument volume is reduced; B6.5.8 requires ineffectiveness to be measured
immediately before rebalancing and treats rebalancing as continuation
\citep[IFRS~9.B6.5.8, B6.5.16--18]{eu20231803}.

Suppose two eligible external derivatives jointly hedge a highly probable
purchase of 100 units, originally designated at price 10. Derivative A has
nominal quantity 60 and B has quantity 40. Both initially cost zero and have
zero value. At the first observation, the purchase price is 11, A is worth 60
and B is worth 80. These complete values and the qualifying relationship are
supplied facts. The total derivative gain is 140, the hedged purchase change
is \(-100\), and the lower-of calculation gives reserve 100 and cumulative
profit 40.

Now suppose a documented change in the relationship requires reducing A's
designated quantity by 30 while the risk-management objective remains
unchanged. Half of A remains designated, all of B remains designated, and the
purchase stays at 100 units. The company still owns both derivatives. The
removed half of A is worth 30 at the reduction date. Taking 30\% of the
aggregate value 140 would give 42, which is wrong: A and B have different
values per nominal unit. IFRS~9.6.2.4(c) permits a proportion of an entire
instrument, and 6.2.5 permits joint designation of combinations
\citep[IFRS~9.6.2.4--6.2.5]{eu20231803}. Those permissions do not make every
instrument unit financially interchangeable.

The reduction itself changes no value, reserve or profit. A's earlier gain of
60 remains in the history, including the 30 attributable to the portion now
removed. Only subsequent changes on that portion are outside the hedge and
enter profit under B6.5.16. At a later price of 12, suppose A is worth 120 and
B is worth 100. The gains accumulated while designated are
\[
 G=140+\tfrac12(120-60)+(100-80)=190.
\]
The removed half of A contributes a further gain
\(U=\tfrac12(120-60)=30\) outside the hedge. Together,
\(G+U=220\), the full value of the two derivatives. The purchase change is
still measured on all 100 units from price 10, giving \(H=-200\).
The reserve is 190 and cumulative profit is
\((190-190)+30=30\). Profit has fallen from 40 to 30 since the previous
observation: reversal of earlier ineffectiveness exceeds the new undesignated
gain by 10.

This history can be represented without treating the remaining nominal ratio
as a valuation rule. For each instrument \(j\), let \(V_j(t)\) be its whole
value, \(w_j(t)\) its remaining designated fraction, and \(a_{jk}\) the
fraction removed at date \(u_{jk}\). All fractions refer to that instrument's
original nominal amount. With zero initial value and no intermediate sale or
settlement, its designated and undesignated gains are
\[
 \begin{aligned}
 G_j(t)&=w_j(t)V_j(t)+\sum_k a_{jk}V_j(u_{jk}),\\
 U_j(t)&=\sum_k a_{jk}\{V_j(t)-V_j(u_{jk})\},\qquad
 w_j(t)+\sum_k a_{jk}=1.
 \end{aligned}
\]
The first line retains each stopped portion's value at the date its contribution
ceased. The second assigns only its later change to profit outside the hedge.
Adding the lines gives \(G_j+U_j=V_j\). An equivalent chronological calculation
adds \(w_j\Delta V_j\) to designated gains and
\((1-w_j)\Delta V_j\) to undesignated gains between observations. Changing
\(w_j\) after a completed measurement creates no jump. This is a local
deduction of prospective rebalancing and cumulative measurement, rather than
an allocation rule for an unidentified reserve.

Sum the component histories to obtain \(G\) and \(U\). While the continuing
hedge has offsetting cumulative instrument and item changes, IFRS~9.6.5.11 gives
\[
 R=\operatorname{sgn}(G)\min(|G|,|H|),\qquad
 \Pi=(G-R)+U.
\]
The reserve remains attached to the unchanged purchase. Deleting a removed
instrument's prior gain would reset history. Applying the earlier matched-
portion method would instead remove purchase units and change the target.
Neither describes this instrument-only reduction. Missing marks, uncertain
qualification, an altered objective or an unidentified portion prevents this
calculation from being admitted.

The identity of the removed instrument matters even when the nominal result
looks identical. Remove 30 units of B instead of A. The stopped fraction is
\(30/40=3/4\), worth 60 at the reduction date. The remaining designated
nominal is again 70. At the same later marks, however,
\[
 G=140+(120-60)+\tfrac14(100-80)=205,\qquad
 U=\tfrac34(100-80)=15.
\]
Reserve is now 200 and profit is \((205-200)+15=20\). The same remaining
nominal ratio therefore produces a different reserve and earnings history.

Complete the original A-reduction case by receiving the 100 units at price 12
and settling both derivatives at the final marks. With opening cash and capital
2000, pay the supplier 1200, receive 120 on A and receive 100 on B. Closing cash
is 1020. Inventory initially costs 1200; transferring reserve 190 into its
basis under IFRS~9.6.5.11(d)(i) leaves inventory 1010
\citep[IFRS~9.6.5.11]{eu20231803}. Cash plus inventory is 2030, equal to capital
2000 plus cumulative profit 30. The supplier payment, two settlements and basis
transfer remain distinct entries. In the B-reduction case, the same cash 1020
joins inventory 1000 and profit 20 instead.

The \texttt{InstrumentReduction/1} report retains component marks, designated
fractions, stopped values, prior gains and settled cash. Eight authored paths
also check unchanged marks, no reduction, proportional component reductions,
a smaller item change, a sign reversal and simultaneous derivative assets and
liabilities. All future changes on owned undesignated portions remain visible.
The method requires a complete purchase-and-settlement bridge and returns no
partial journals on refusal. Supplied valuation, recoverability, qualifying
facts and documentation remain substantive prerequisites. Actual and aligned
option time values, tax, cash-flow activity classification and allocation among
heterogeneous discontinued purchases require their separate methods.
